\begin{figure}[htbp]
  \centering
  \begin{tikzpicture}[alt={Two overlapping eight-sample windows for a single impulse. The first seven samples are zero history. Original sample seven is 16000, followed by zeros. Output zero uses indices zero through seven; output one uses one through eight. The impulse remains in each window through output seven; output eight no longer contains it. With every coefficient one eighth, outputs zero through seven are 2000 and output eight is zero.},x=.67cm,y=.65cm]
    \foreach \j in {0,...,9} {
        \draw (\j,0) rectangle +(1,1);
        \node[font=\scriptsize] at (\j+.5,1.35) {\j};
      }
    \foreach \j in {0,...,6,8,9} {\node[font=\small] at (\j+.5,.5) {0};}
    \node[font=\scriptsize\bfseries] at (7.5,.5) {16000};
    \draw[line width=1pt] (0,-.25)--(0,-.5)--(8,-.5)--(8,-.25);
    \node[anchor=east,font=\small] at (-.15,-.5) {$i=0$};
    \draw[line width=1pt] (1,-.85)--(1,-1.1)--(9,-1.1)--(9,-.85);
    \node[anchor=east,font=\small] at (.85,-1.1) {$i=1$};
    \node[anchor=west,font=\small] at (8.2,-.5) {$2000$};
    \node[anchor=west,font=\small] at (9.2,-1.1) {$2000$};
  \end{tikzpicture}
  \caption{A window moves one sample for each output. Seven history samples make the first complete window available. The impulse initially meets coefficient seven, then coefficient six, and so on until it leaves the window.}
  \label{fig:revision-fir-windows}
\end{figure}
